% Generated draft: requires mathematical and source review.
\documentclass[11pt]{article}
\usepackage[margin=25mm]{geometry}
\usepackage{fontspec}
\setmainfont{Latin Modern Roman}
\usepackage{amsmath,amssymb,mathtools,unicode-math}
\setmathfont{Latin Modern Math}
\usepackage{enumitem,needspace,xurl,hyperref}
\setlength{\parindent}{0pt}
\setlength{\parskip}{5pt}
\setlength{\emergencystretch}{3em}
\newcommand{\N}{\mathbb{N}}
\newcommand{\Z}{\mathbb{Z}}
\newcommand{\Q}{\mathbb{Q}}
\newcommand{\R}{\mathbb{R}}
\newcommand{\C}{\mathbb{C}}
\newcommand{\F}{\mathbb{F}}
\newcommand{\A}{\mathbb{A}}
\newcommand{\PP}{\mathbb{P}}
\newcommand{\E}{\mathbb{E}}
\newcommand{\eps}{\varepsilon}
\newcommand{\dd}{\,\mathrm{d}}
\begin{document}
\section*{UnsolvedMath ID: 11000102}
\textbf{Source problem code:} catalogue entry
\subsection{Definitions and mathematical statement}
Let \(d_W\) be the word metric for a fixed finite generating set \(W\) of the mapping class group in the source, and let \(t\) be its specified element. The preceding definition of \(t\), surface and generators is absent. The question asks for the growth of
\[n\longmapsto d_W(t^n,1)\qquad(n\to\infty).\]
For an infinite-order element, linear growth means that there are constants \(a>0,b\geq0,c>0\), independent of \(n\), with
\[an-b\leq d_W(t^n,1)\leq cn\qquad(n\geq1).\]
The upper linear bound follows from the word metric; the substantive assertion is a positive linear lower bound, equivalently undistortion of the cyclic subgroup in this setting. If \(t\) were finite order its powers would be bounded, so the intended element hypothesis is essential. The extraction's \(tn\) is read as \(t^n\), consistently with the stated growth problem, not as a product with an integer element.

The source says one might expect either linear or \(O(\log n)\) behavior, gives the higher-rank arithmetic-group comparison, and then expressly reports that Farb--Lubotzky--Minsky [FLM] proved linear growth, solving this question completely. This solved qualification is retained; no additional open dichotomy is asserted. The source attributes logarithmic growth for virtually unipotent elements in the indicated arithmetic rank-\(\geq2\) setting to Lubotzky--Moses--Raghunathan [LMR]. Its exact arithmetic hypotheses and the scope of \(t\) in the mapping-class result require the missing references. The comparison is not an assertion that the two settings have identical hypotheses.
\subsection{Short English statement}
The source asks how word length grows for powers of the specified mapping class, and then states that Farb--Lubotzky--Minsky proved linear growth. The element and surface hypotheses must be recovered from context.
\subsection{Sources}
\begin{itemize}
\item Local UnsolvedMath registry, permanent ID \(11000102\), exact source code \texttt{catalogue entry}.
\item Farb, \emph{Problems on Mapping Class Groups and Related Topics}, unnumbered Question, PDF page 85: \url{https://www.math.uchicago.edu/~farb/papers/mcgbook.pdf}. Registry access 2026-07-29; pdf-font-structure extraction. The running text identifies the chapter as \emph{Fifteen problems about the mapping class groups}.
\item B. Farb, A. Lubotzky and Y. Minsky [FLM] are credited with the stated linear-growth resolution. A. Lubotzky, S. Moses and M. S. Raghunathan [LMR] are cited for the arithmetic comparison. Full bibliographic entries are absent.
\end{itemize}
\end{document}
